% Source: Weinberg GR, physical PDF pp. 92--95
%         (printed pp. 67--70).
% Coverage: Section 3.1, Statement of the Principle; equations
%           (3.1.1)--(3.1.3).
% Figures/tables/footnotes: none; reference marker 1.
% Status: source-reviewed and compile-clean.
% Uncertainties: none.

\subsection{Statement of the Principle}\label{sec:3.1}

The Principle of Equivalence rests on the equality of gravitational and
inertial mass, demonstrated by Galileo, Huygens, Newton, Bessel, and
E\"otv\"os. (See Section~1.2.) Einstein reflected that, as a consequence, no
external static homogeneous gravitational field could be detected in a
freely falling elevator, for the observers, their test bodies, and the
elevator itself would respond to the field with the same acceleration. This
can be easily proved for a system of particles $N$, moving with
nonrelativistic velocities under the influence of forces
$\mathbf F(\mathbf x_N-\mathbf x_M)$ (e.g., electrostatic or gravitational
forces) and an external gravitational field $\mathbf g$. The equations of
motion are
\begin{equation}
  m_N\frac{\dd^2\mathbf x_N}{\dd t^2}
  =m_N\mathbf g+\sum_M\mathbf F(\mathbf x_N-\mathbf x_M).
  \tag{3.1.1}\label{eq:3.1.1}
\end{equation}
Suppose that we perform a non-Galilean spacetime coordinate transformation
\begin{equation}
  \mathbf x'=\mathbf x-\frac12\mathbf g t^2,
  \qquad
  t'=t.
  \tag{3.1.2}\label{eq:3.1.2}
\end{equation}
Then $\mathbf g$ will be canceled by an inertial ``force,'' and the equation
of motion will become
\begin{equation}
  m_N\frac{\dd^2\mathbf x'_N}{\dd t'^2}
  =\sum_M\mathbf F(\mathbf x'_N-\mathbf x'_M).
  \tag{3.1.3}\label{eq:3.1.3}
\end{equation}
Hence the original observer $O$, who uses coordinates $\mathbf x,t$, and his
freely falling friend $O'$, who uses $\mathbf x',t'$, will detect no
difference in the laws of mechanics, except that $O$ will say that he feels a
gravitational field and $O'$ will say that he does not. The equivalence
principle says that this cancellation of gravitational by inertial force (and
hence their equivalence) will obtain for \emph{all} freely falling systems,
whether or not they can be described by simple equations such as
\eqref{eq:3.1.1}.

We are not yet ready to state the Principle of Equivalence in its final form,
because the preceding remarks dealt only with a static homogeneous
gravitational field. Had $\mathbf g$ depended on $\mathbf x$ or $t$, we would
not have been able to eliminate it from the equations of motion by the
acceleration \eqref{eq:3.1.2}. For example, the earth is in free fall about
the sun, and for the most part we on earth do not feel the sun's
gravitational field, but the slight inhomogeneity in this field (about 1 part
in 6000 from noon to midnight) is enough to raise impressive tides in our
oceans. Even the observers in Einstein's freely falling elevator would in
principle be able to detect the earth's field, because objects in the
elevator would be falling radially toward the center of the earth, and hence
would approach each other as the elevator descended.

Although inertial forces do not exactly cancel gravitational forces for
freely falling systems in an inhomogeneous or time-dependent gravitational
field, we can still expect an approximate cancellation if we restrict our
attention to such a small region of space and time that the field changes
very little over the region. Therefore we formulate the equivalence principle
as the statement that \emph{at every spacetime point in an arbitrary
gravitational field it is possible to choose a ``locally inertial coordinate
system'' such that, within a sufficiently small region of the point in
question, the laws of nature take the same form as in unaccelerated Cartesian
coordinate systems in the absence of gravitation}. There is a little
vagueness here about what we mean by ``the same form as in unaccelerated
Cartesian coordinate systems,'' so to avoid any possible ambiguity we can
specify that by this we mean the form given to the laws of nature by special
relativity, for example, such equations as (2.3.1), (2.7.6), (2.7.7),
(2.7.9), and (2.8.7). There is also a question of how small is ``sufficiently
small.'' Roughly speaking, we mean that the region must be small enough so
that the gravitational field is sensibly constant throughout it, but we
cannot be more precise until we learn how to represent the gravitational
field mathematically. (See the end of Section~4.1.)

The attentive reader may have noticed a certain resemblance between the
Principle of Equivalence and the axiom which Gauss took as the basis of
non-Euclidean geometry. The Principle of Equivalence says that at any point
in spacetime we may erect a locally inertial coordinate system in which
matter satisfies the laws of special relativity. We saw in Chapter~1 that
Gauss assumed that at any point on a curved surface we may erect a locally
Cartesian coordinate system in which distances obey the law of Pythagoras.
Because of this deep analogy, we should expect the laws of gravitation to
bear a strong resemblance to the formulas of Riemannian geometry. In
particular, Gauss's assumption implies that all inner properties of a curved
surface can be described in terms of derivatives
$\partial\xi^\rho/\partial x^\mu$ of the function $\xi^\rho(x)$ that defines
the transformation $x\to\xi$ from some general coordinate system $x^\mu$
covering the surface to the locally Cartesian system $\xi^\rho$, whereas the
Principle of Equivalence tells us that all effects of a gravitational field
can be described in terms of derivatives
$\partial\xi^\rho/\partial x^\mu$ of the function $\xi^\rho(x)$ that defines
the transformation from the ``laboratory'' coordinates $x^\mu$ to the
locally inertial coordinates $\xi^\rho$. Furthermore, it was shown in
Chapter~1 that the geometrically relevant functions of these derivatives are
the quantities $g_{\mu\nu}$ defined by Eq.~(1.1.7); we shall see in the
following sections of this chapter that the gravitational field is described
in just the same way.

Occasionally one finds references to a ``weak Principle of Equivalence'' and
a ``strong Principle of Equivalence.'' The strong Principle of Equivalence is
just what I have already stated, with ``laws of nature'' meaning \emph{all}
the laws of nature. The weak principle is the same, but with ``laws of
nature'' replaced by ``laws of motion of freely falling particles.'' That is,
the weak principle is nothing but a restatement of the observed equality of
gravitational and inertial mass, whereas the strong principle is a
generalization of these observations that governs the effects of gravitation
on all physical systems.

The experiments of E\"otv\"os, Dicke, and their predecessors (see
Section~1.2) provide direct verification only of the weak Principle of
Equivalence, but they provide some indirect evidence for the strong
principle. The mass of different substances arises \emph{in different
proportions} from the masses of the neutrons and protons plus electrons of
which they are composed, and from the strong and electromagnetic forces that
bind these particles together, so the ratio of gravitational to inertial
mass will be equal for all these substances only if it is equal for their
constituents. Wapstra and Nijgh~\hyperref[ch3-ref-1]{[1]} have shown that the
limits set by E\"otv\"os on any possible inequality in the ratio of
gravitational to inertial mass for glass, cork, antimonite, and brass imply
that this ratio is equal for neutrons and protons plus electrons to 1 part in
$6\times10^5$ and equal for neutrons and binding energies to 1 part in
$1.2\times10^4$. To this accuracy, an observer in a freely falling
coordinate system will detect no gravitational force on neutrons, hydrogen,
or their binding energies. It would be difficult to conceive of a theory
that satisfies this requirement and does not go all the way to the strong
principle (that no gravitational effects of any sort can be felt in a
locally inertial frame).

We might, however, distinguish two versions of the strong principle of
equivalence, a ``very strong principle,'' which applies to all phenomena, and
a ``medium-strong principle,'' which applies to all phenomena except
gravitation itself. Certainly the experiments of E\"otv\"os and Dicke are not
accurate enough to say whether gravitational binding energies affect
inertial and gravitational masses in the same way. This question might be
settled by studying the motion of a small body in orbit about a large body
that is itself in free fall in a gravitational field. For instance, the
gravitational binding energy of the earth contributes a fraction
$-8.4\times10^{-10}$ of its total mass, whereas the gravitational binding
energy of an artificial satellite contributes a very much smaller fraction
of its mass. Thus, if (to take an extreme case) the (negative) gravitational
binding energy contributes fully to the inertial mass but not at all to the
gravitational mass, then the ratio of gravitational to inertial mass of the
satellite would be greater than that for the earth by a fraction
$8.4\times10^{-10}$. The earth is in free fall, with the gravitational
attraction of the sun balanced by the inertial force owing to the earth's
revolution. The gravitational and inertial forces on the satellite owing to
the presence of the sun and the earth's revolution are equal (neglecting for
a moment the distance between the satellite and the earth's center of mass)
to the gravitational and inertial forces on the earth times the ratio of
gravitational or inertial masses, so these two forces are not in balance for
the satellite, the gravitational force being greater than the inertial force
by a fraction $8.4\times10^{-10}$. The acceleration owing to the sun's
gravity is at the orbit of the earth about $6\times10^{-4}$ of the
acceleration owing to the earth's gravity at the surface of the earth, so we
conclude that if the gravitational binding energy of the earth contributed
fully to its inertial mass but not at all to its gravitational mass, then an
artificial satellite in a low orbit about the earth would feel an effective
attraction toward the sun equal to about $5.4\times10^{-13}$ times its
gravitational attraction toward the earth. This tiny effect would be
entirely masked by a ``tidal'' force because the satellite is far from the
center of mass of the earth, and there is no prospect of its being measured.
This is a pity, because it is precisely the very strong assumption that the
Principle of Equivalence applies to gravitational fields that will lead us
in Chapter~5 to Einstein's field equations for gravitation.
